% GENERATED FILE. Edit canonical lessons, then run npm run generate:books.
% canonical-source-hash: fnv1a64:b469758643169f1c
% canonical-lessons: ZH-C117-dongbian, ZH-C117-nanbian, ZH-C117-xibian, ZH-C117-beibian, ZH-C117-fangxiang

\chapter{Which Way?}
\label{ch:zh-which-way}
\hlchaptermodality{117}

\begin{chapteropening}
\textbf{By the end of this chapter:} \emph{I can say east, south, west and north, and ask which way to go.}

Complete the last lesson of chapter 117: I can say east, south, west and north, and ask which way to go.
\end{chapteropening}

\section[dōngbian]{dōngbian — the east; on the east side}

\label{lesson:ZH-C117-dongbian}



\begin{quote}
Before the new one: say the Mandarin for umbrella, then the Mandarin for to lose.
\end{quote}

\subsection*{You'll want to know: dōngbian}
\emph{dōngbian} — \textquotedblleft{}the east; on the east side\textquotedblright{}.

\begin{quote}
\emph{chēzhàn zài dōngbian} — The station is on the east side.
\end{quote}

Its characters come later; for now it is a word for the ear.

\begin{cousinweb}
\emph{dōngbian} is built from dōng 'east' and \zh{边} 'side', both inherited from Middle Chinese, \emph{tuwng} and \emph{pen} (Guangyun, in Baxter's transcription); the dictionary calls the east the direction of spring, and repeats an old reading of the character as the sun inside a tree.
\end{cousinweb}

\begin{grammarlens}[title={What you've built}]
One more word for asking where things are.
\end{grammarlens}

\subsection*{Your turn}
\begin{itemize}
\raggedright
  \item \emph{Say it:} \emph{dōngbian}
  \item \emph{Say it:} \emph{dōngbian}, once more
  \item \emph{Say it:} \emph{chēzhàn zài dōngbian} — The station is on the east side.
  \item \emph{Recall:} say the Mandarin for umbrella, then the Mandarin for to lose, then say \emph{dōngbian} again
\end{itemize}

\begin{tcolorbox}[breakable,colback=teal!4,colframe=teal!35!black,title={Before you move on}]
What does dōngbian mean? (\textquotedblleft{}The east; on the east side\textquotedblright{}.) Say it once more.
\end{tcolorbox}
\section[nánbian]{nánbian — the south; on the south side}

\label{lesson:ZH-C117-nanbian}



\begin{quote}
Before the new one: say the Mandarin for to lose, then the Mandarin for the east.
\end{quote}

\subsection*{You'll want to know: nánbian}
\emph{nánbian} — \textquotedblleft{}the south; on the south side\textquotedblright{}.

\begin{quote}
\emph{gōngyuán zài nánbian} — The park is to the south.
\end{quote}

Its characters come later; for now it is a word for the ear.

\begin{cousinweb}
\emph{nánbian} is built from nán 'south' and \zh{边} 'side', both inherited from Middle Chinese, \emph{nom} and \emph{pen} (Guangyun, in Baxter's transcription); the dictionary calls the south the direction of fire.
\end{cousinweb}

\begin{grammarlens}[title={What you've built}]
One more word for asking where things are.
\end{grammarlens}

\subsection*{Your turn}
\begin{itemize}
\raggedright
  \item \emph{Say it:} \emph{nánbian}
  \item \emph{Say it:} \emph{nánbian}, once more
  \item \emph{Say it:} \emph{gōngyuán zài nánbian} — The park is to the south.
  \item \emph{Recall:} say the Mandarin for to lose, then the Mandarin for the east, then say \emph{nánbian} again
\end{itemize}

\begin{tcolorbox}[breakable,colback=teal!4,colframe=teal!35!black,title={Before you move on}]
What does nánbian mean? (\textquotedblleft{}The south; on the south side\textquotedblright{}.) Say it once more.
\end{tcolorbox}
\section[xībian]{xībian — the west; on the west side}

\label{lesson:ZH-C117-xibian}



\begin{quote}
Before the new one: say the Mandarin for the east, then the Mandarin for the south.
\end{quote}

\subsection*{You'll want to know: xībian}
\emph{xībian} — \textquotedblleft{}the west; on the west side\textquotedblright{}.

\begin{quote}
\emph{xībian yǒu shān} — There are hills to the west.
\end{quote}

Its characters come later; for now it is a word for the ear.

\begin{cousinweb}
\emph{xībian} is built from xī 'west' and \zh{边} 'side', both inherited from Middle Chinese, \emph{sej} and \emph{pen} (Guangyun, in Baxter's transcription); the dictionary calls the west the direction of autumn, and explains the character as a bird settling on its nest when the sun is in the west.
\end{cousinweb}

\begin{grammarlens}[title={What you've built}]
One more word for asking where things are.
\end{grammarlens}

\subsection*{Your turn}
\begin{itemize}
\raggedright
  \item \emph{Say it:} \emph{xībian}
  \item \emph{Say it:} \emph{xībian}, once more
  \item \emph{Say it:} \emph{xībian yǒu shān} — There are hills to the west.
  \item \emph{Recall:} say the Mandarin for the east, then the Mandarin for the south, then say \emph{xībian} again
\end{itemize}

\begin{tcolorbox}[breakable,colback=teal!4,colframe=teal!35!black,title={Before you move on}]
What does xībian mean? (\textquotedblleft{}The west; on the west side\textquotedblright{}.) Say it once more.
\end{tcolorbox}
\section[běibian]{běibian — the north; on the north side}

\label{lesson:ZH-C117-beibian}



\begin{quote}
Before the new one: say the Mandarin for the south, then the Mandarin for the west.
\end{quote}

\subsection*{You'll want to know: běibian}
\emph{běibian} — \textquotedblleft{}the north; on the north side\textquotedblright{}.

\begin{quote}
\emph{Běijīng zài Zhōngguó de běibian} — Beijing is in the north of China.
\end{quote}

Its characters come later; for now it is a word for the ear.

\begin{cousinweb}
\emph{běibian} is built from běi 'north' and \zh{边} 'side', both inherited from Middle Chinese, \emph{pok} and \emph{pen} (Guangyun, in Baxter's transcription).
\end{cousinweb}

\begin{grammarlens}[title={What you've built}]
One more word for asking where things are.
\end{grammarlens}

\subsection*{Your turn}
\begin{itemize}
\raggedright
  \item \emph{Say it:} \emph{běibian}
  \item \emph{Say it:} \emph{běibian}, once more
  \item \emph{Say it:} \emph{Běijīng zài Zhōngguó de běibian} — Beijing is in the north of China.
  \item \emph{Recall:} say the Mandarin for the south, then the Mandarin for the west, then say \emph{běibian} again
\end{itemize}

\begin{tcolorbox}[breakable,colback=teal!4,colframe=teal!35!black,title={Before you move on}]
What does běibian mean? (\textquotedblleft{}The north; on the north side\textquotedblright{}.) Say it once more.
\end{tcolorbox}
\section[fāngxiàng]{fāngxiàng — direction; way}

\label{lesson:ZH-C117-fangxiang}



\begin{quote}
Before the new one: say the Mandarin for the west, then the Mandarin for the north.
\end{quote}

\subsection*{You'll want to know: fāngxiàng}
\emph{fāngxiàng} — \textquotedblleft{}direction; way\textquotedblright{}.

\begin{quote}
\emph{chēzhàn zài nǎge fāngxiàng} — Which way is the station?
\end{quote}

Its characters come later; for now it is a word for the ear.

\begin{cousinweb}
\emph{fāngxiàng} is built from fāng 'side, region' and xiàng 'to face', both inherited from Middle Chinese, \emph{pjang} and \emph{xjangH} (Guangyun, in Baxter's transcription): the side you face.
\end{cousinweb}

\begin{grammarlens}[title={What you've built}]
One more word for asking where things are.
\end{grammarlens}

\subsection*{Your turn}
\begin{itemize}
\raggedright
  \item \emph{Say it:} \emph{fāngxiàng}
  \item \emph{Say it:} \emph{fāngxiàng}, once more
  \item \emph{Say it:} \emph{chēzhàn zài nǎge fāngxiàng} — Which way is the station?
  \item \emph{Recall:} say the Mandarin for the west, then the Mandarin for the north, then say \emph{fāngxiàng} again
\end{itemize}

\begin{tcolorbox}[breakable,colback=teal!4,colframe=teal!35!black,title={Before you move on}]
What does fāngxiàng mean? (\textquotedblleft{}Direction; way\textquotedblright{}.) Say it once more.
\end{tcolorbox}
